
\chapter{Transaction API Design – Software Engineering Perspectives}
\section{Classic Synchronous API}
\begin{itemize}
    \item \texttt{tm\_begin()}, \texttt{tm\_read()}, \texttt{tm\_write()}, \texttt{tm\_commit()}, \texttt{tm\_abort()}.
    \item Thread‑local transaction descriptors and the need for thread‑entry instrumentation.
    \item Typed wrappers (\texttt{tm\_read\_i4}, \texttt{tm\_write\_i4}, \ldots) and how they map onto the byte-oriented core.
    \item The retry loop: \texttt{sigsetjmp}/\texttt{siglongjmp} (or equivalent) to restart aborted transactions --- and the frame-lifetime hazard this creates (Part~VI).
    \item The set of primitives is deliberately small; everything else (nested transactions, retry, allocation) is built on top.
\end{itemize}

\begin{figure}[htbp]
\centering
\resizebox{\textwidth}{!}{%
\begin{tikzpicture}[
    node distance=1.8cm and 2.4cm,
    state/.style={draw, rounded corners, minimum width=2.8cm, minimum height=0.9cm, align=center},
    arrow/.style={-{Stealth}, thick}
]
\node[state] (outside) {Outside\\transaction};
\node[state, right=of outside] (active) {Active\\read/write set};
\node[state, right=of active] (commit) {Commit\\validate};
\node[state, above right=1.2cm and 1.7cm of commit] (done) {Committed\\visible};
\node[state, below right=1.2cm and 1.7cm of commit] (retry) {Aborted\\retry};

\draw[arrow] (outside) -- node[above]{\texttt{tm\_begin}} (active);
\draw[arrow] (active) edge[loop above] node{\texttt{tm\_read}/\texttt{tm\_write}} (active);
\draw[arrow] (active) -- node[above]{\texttt{tm\_end}} (commit);
\draw[arrow] (commit) -- node[above left]{valid} (done);
\draw[arrow] (commit) -- node[below left]{conflict} (retry);
\draw[arrow] (retry) |- node[pos=0.25, below]{loop} (outside);
\end{tikzpicture}%
}
\caption{Control states exposed by the synchronous transaction API. Backends differ in validation and metadata, but the source-level protocol is the same.}
\label{fig:api-state-machine}
\end{figure}

\section{Worked Example: The Transfer, on the Classic API}
\index{TMD API}\index{tm\_begin}\index{tm\_end}\index{retry loop}
\label{sec:api-example}

Chapter~1's running example, expressed with the synchronous API and an
explicit retry loop:

\begin{lstlisting}[language=C, caption={transfer on the synchronous API.}]
int transfer(int *a, int *b, int n) {
    int attempts = 0;
    for (;;) {
        tm_begin();
        attempts++;
        int ba = tm_read_i4(a);
        int bb = tm_read_i4(b);
        if (ba - n < 0 || bb + n < 0) { tm_end(); return 0; }
        tm_write_i4(a, ba - n);
        tm_write_i4(b, bb + n);
        if (tm_end() == TM_COMMITTED) return attempts;  /* done */
        /* aborted: loop and retry */
    }
}
\end{lstlisting}

\begin{itemize}
    \item Compare with Chapter~1's lock-based version: no lock ordering, no unlock on every early-return path, and the two writes are atomic by construction.
    \item The \texttt{sigsetjmp}/\texttt{siglongjmp} machinery hides in \texttt{tm\_begin} and \texttt{tm\_end} --- the frame-lifetime hazard of Part~VI.
    \item Every backend in Part~IV implements exactly these five functions, so this source compiles unchanged across all of them.
    \item Exercise~3 asks for the \texttt{\_\_transaction\_atomic} equivalent; note the retry loop disappears --- the compiler generates it.
\end{itemize}

\section{Worked Example: A Lambda Wrapper}
\index{transactional lambda}\index{C++ wrapper}\index{retry loop}

A thin C++ wrapper can remove duplicated retry loops while preserving the same
backend contract. The body is still responsible for using instrumented loads and
stores.

\begin{lstlisting}[language=C++, caption={a C++ retry wrapper over the classic API.}]
extern "C" {
enum { TM_ABORTED = 0, TM_COMMITTED = 1 };

void tm_begin(void);
int  tm_end(void);
int  tm_read_i4(const int *addr);
void tm_write_i4(int *addr, int value);
}

template <class Body>
auto atomic_retry(Body body) -> decltype(body()) {
    for (;;) {
        tm_begin();

        auto result = body();

        if (tm_end() == TM_COMMITTED)
            return result;

        /* aborted: discard result and execute body again */
    }
}

int transfer_with_floor(int *from, int *to, int amount, int from_floor) {
    return atomic_retry([=]() -> int {
        if (amount <= 0)
            return 0;

        int bf = tm_read_i4(from);
        int bt = tm_read_i4(to);

        if (bf - amount < from_floor)
            return 0;

        tm_write_i4(from, bf - amount);
        tm_write_i4(to,   bt + amount);
        return 1;
    });
}
\end{lstlisting}

The money-conservation step for the successful path is local to the closure:

$$
\mathit{from}' + \mathit{to}' =
(\mathit{from} - \mathit{amount}) + (\mathit{to} + \mathit{amount})
=
\mathit{from} + \mathit{to}
$$

\begin{itemize}
    \item The wrapper centralises retry policy; the transfer code only describes the atomic update.
    \item The closure must be replay-safe. Printing, freeing memory, or sending messages inside the body would repeat on abort.
    \item This is not language integration. Ordinary loads and stores are still ordinary loads and stores; the wrapper cannot instrument them.
    \item Destructors are not a complete substitute for the API. If a backend restarts with \texttt{siglongjmp}, C++ stack unwinding does not run.
\end{itemize}

\section{Language‑Integrated Keywords}
\begin{itemize}
    \item GCC’s \texttt{\_\_transaction\_atomic} and its implicit instrumentation.
    \item How keywords can simplify the programmer’s view.
    \item Comparison with explicit API: the compiler inserts begin/commit and instruments accesses automatically, removing the risk of a forgotten \texttt{tm\_end}.
    \item The language-level cost: keywords need compiler support for exception handling, cancellation, and safe/unsafe functions (see Chapter~4 on the C++ TS).
\end{itemize}

\begin{table}[htbp]
\centering
\begin{tabular}{p{0.19\linewidth}p{0.23\linewidth}p{0.23\linewidth}p{0.23\linewidth}}
\hline
API style & Instrumentation point & Retry owner & Main engineering risk \\
\hline
Classic calls & Programmer / wrapper & Application loop & Forgotten access or end path \\
Language keyword & Compiler & Compiler/runtime & Compiler and ABI support \\
Typed object wrapper & Library type & Library/runtime & Escaping raw references \\
Future-based closure & Executor worker & Executor & Replay of side effects \\
Distributed transaction & Client coordinator & Coordinator & Partial failure and timeouts \\
\hline
\end{tabular}
\caption{API families used in TM systems. The lower-level runtime may still be TL2, NOrec, SGL, Percolator-style two-phase commit, or another backend; the table classifies the programming interface.}
\label{tab:api-taxonomy}
\end{table}

\section{Asynchronous and Future‑Based APIs}
\begin{itemize}
    \item Pushing transactions into work queues.
    \item Transactional futures: submitting a transaction as a future, awaiting results.
    \item Batching and potential for optimisation.
    \item The executor model: worker threads consume transaction closures from a queue; transactions become units of work, not thread-stack entities.
    \item Pipeline implications: with async execution, stack-address checks and thread-entry initialization become per-task concerns, not per-thread.
    \item Composing futures preserves the composability argument of Chapter~1: an \texttt{atomic} closure is just another value.
\end{itemize}

\section{Worked Example: Futures for Batched Transfers}
\index{transactional future}\index{executor}\index{batching}

A future-based API moves the retry loop into the worker. The caller receives a
handle and later observes the committed result.

\begin{lstlisting}[language=C++, caption={submitting bank transfers as transactional futures.}]
#include <future>
#include <utility>
#include <vector>

extern "C" {
enum { TM_ABORTED = 0, TM_COMMITTED = 1 };

void tm_begin(void);
int  tm_end(void);
int  tm_read_i4(const int *addr);
void tm_write_i4(int *addr, int value);
}

/* Concrete helper for transactions returning int. */
template <class Body>
std::future<int> submit_tx_int(Body body) {
    return std::async(std::launch::async,
        [body = std::move(body)]() mutable -> int {
            for (;;) {
                tm_begin();

                int result = body();

                if (tm_end() == TM_COMMITTED)
                    return result;

                /* aborted: retry on this worker */
            }
        });
}

std::future<int> submit_transfer(int *from, int *to,
                                 int amount, int from_floor) {
    return submit_tx_int([=]() -> int {
        if (amount <= 0)
            return 0;

        int bf = tm_read_i4(from);
        int bt = tm_read_i4(to);

        if (bf - amount < from_floor)
            return 0;

        tm_write_i4(from, bf - amount);
        tm_write_i4(to,   bt + amount);
        return 1;
    });
}

int submit_batch_and_count(int *a, int *b, int *c) {
    std::vector<std::future<int>> fs;

    fs.push_back(submit_transfer(a, b, 10, 0));
    fs.push_back(submit_transfer(b, c,  7, 0));
    fs.push_back(submit_transfer(c, a,  3, 0));

    int committed_successes = 0;
    for (auto &f : fs)
        committed_successes += f.get();

    return committed_successes;
}
\end{lstlisting}

\begin{itemize}
    \item The caller no longer owns transaction stack state; the worker does.
    \item Futures make composition explicit: submit many transactions, then join at the dependency point.
    \item A real executor can batch validation or scheduling decisions. The source-level API does not require the caller to know which backend is underneath.
    \item The closure must capture stable addresses. Capturing pointers to caller stack variables is unsafe if the future may outlive the stack frame.
\end{itemize}

\section{Object‑Oriented and Polymorphic Interfaces}
\begin{itemize}
    \item Template-based STMs (e.g., RSTM).
    \item Trade-offs: expressiveness vs.\ runtime overhead.
    \item Type-safe wrappers such as \texttt{TM<int>} hide read\slash{}write instrumentation behind \texttt{operator*}.
    \item Visitor/transaction-object patterns: objects that register their own read/write sets with the runtime.
    \item The object-granularity problem: per-object locking is coarse; word-level STM reclaims granularity at the cost of metadata indexing.
\end{itemize}

\section{Summary}
\begin{itemize}
    \item API design directly affects usability, safety, and performance.
    \item Asynchronous interfaces open the door to batching and higher throughput at the cost of programming complexity.
    \item The API is the contract the rest of the book's runtimes implement: every backend in Part~IV exposes the same synchronous primitives.
\end{itemize}

\section{Exercises}
\begin{enumerate}
    \item A programmer argues that \texttt{tm\_write} should return a status code indicating whether the write caused an immediate abort. Explain why this violates opacity and propose an alternative.
    \item Sketch an asynchronous API where multiple transactions are submitted to a work queue and later awaited. Discuss how batching might reduce validation overhead.
    \item Compare using \texttt{tm\_begin}\slash{}\texttt{tm\_end} versus \texttt{\_\_transaction\_\allowbreak atomic}. Which is less error-prone and why?
    \item \emph{[implementation]} Extend Listing~\ref{sec:api-example}'s transfer into a reusable \texttt{atomic\_retry} helper for transactions returning \texttt{int}. Add a test that forces at least one abort and checks that the final account sum is unchanged.
    \item \emph{[verification]} Write a small PlusCal or TLA+ model of the state machine in Figure~\ref{fig:api-state-machine}. Check that no execution reaches \emph{Committed} after a failed validation edge.
    \item \emph{[theory]} Compare explicit calls, compiler keywords, and future-based closures under opacity. Which API makes it easiest to prevent non-transactional side effects from becoming visible before commit?
\end{enumerate}
